\section{Valuations Beyond Frames}
\label{sec:lattice}

Definition~\ref{def:frame} fixes the event algebra to be a frame: a
complete lattice in which finite meets distribute over arbitrary
joins $\bigsqcup_{i \in I} a_i$, the operation Heyting implication
$\Rightarrow$ (the right adjoint to $a \sqcap (-)$) is built from.
Definition~\ref{def:valuation}'s four axioms, normalization at
$\bot$ and $\top$, monotonicity, and modularity
$v(a) + v(b) = v(a \sqcup b) + v(a \sqcap b)$, use only the
\emph{binary} join $\sqcup$; none of the four mentions the
\emph{arbitrary} join $\bigsqcup_{i \in I}$. Completeness is used
only where later sections evaluate a valuation on an infinite join,
as in the Scott-continuity results of
Section~\ref{sec:representation}, or build $\Rightarrow$ and $\neg$
from it, as in Section~\ref{sec:hinge}'s slack.
Section~\ref{sec:conclusion} leaves open how much of the
construction is tied to frames specifically, and whether some
ambient category sends an event algebra to a probability calculus
in general. This section settles both questions for one fragment of
the construction, the sharp/point correspondence of
Section~\ref{sec:hinge}: it isolates the exact algebraic hypothesis
that correspondence needs and shows the hypothesis cannot be
dropped. It does not recover $\Rightarrow$ or $\neg$ without
completeness: the adjoint defining $\Rightarrow$ is built from an
arbitrary join of witnesses, so it is unavailable once completeness
goes, even if only finite joins are missing. The valuations studied
below are negation-free by construction, not frame valuations with
completeness quietly removed.

\paragraph{The functor without completeness.}
A \emph{bounded lattice} is a lattice with a least element $\bot$
and a greatest element $\top$; a \emph{bounded lattice homomorphism}
preserves $\bot$, $\top$, $\sqcup$, and $\sqcap$.
Definition~\ref{def:valuation}'s axioms make sense over any bounded
lattice $\Omega$, complete or not:

\begin{definition}[lattice valuation]
\label{def:lvaluation}
A \emph{lattice valuation} on a bounded lattice $\Omega$ is a
monotone function $v : \Omega \to \ENNReal$ with $v(\bot) = 0$,
$v(\top) = 1$, and, for all $a, b \in \Omega$,
$v(a) + v(b) = v(a \sqcup b) + v(a \sqcap b)$.
\end{definition}

\begin{lstlisting}
structure LValuation (Ω : Type*) [Lattice Ω] [BoundedOrder Ω] where
  toFun : Ω → ℝ≥0∞
  map_bot' : toFun ⊥ = 0
  map_top' : toFun ⊤ = 1
  mono'    : Monotone toFun
  modular' : ∀ a b, toFun a + toFun b
             = toFun (a ⊔ b) + toFun (a ⊓ b)
\end{lstlisting}

\noindent
A lattice valuation on a frame is exactly a valuation in the sense
of Definition~\ref{def:valuation}; the two notions differ only in
how much structure the ambient $\Omega$ carries.

Section~\ref{sec:monad} pulls a valuation back along a frame
homomorphism and checks the assignment is functorial (identity,
composition). That argument pattern-matches on $\sqcup$, $\sqcap$,
$\bot$, $\top$ alone and never touches $\bigsqcup$:

\begin{proposition}[pullback without completeness]
\label{prop:lvaluation-comap}
For bounded lattices $\Gamma, \Omega$, a bounded lattice
homomorphism $f : \Gamma \to \Omega$, and a lattice valuation $v$ on
$\Omega$, $v \circ f$ is a lattice valuation on $\Gamma$. This
assignment sends the identity homomorphism to the identity and
composition of homomorphisms to composition of the induced maps.
\end{proposition}

\noindent
(Mechanized as \texttt{LValuation.comap}, \texttt{comap\_id}, and
\texttt{comap\_comp} in \texttt{GeneralValuation.\allowbreak lean}.)
$\mathrm{Val}(-)$ is therefore a functor on bounded lattices with
bounded lattice homomorphisms, of which the frame-level functor of
Section~\ref{sec:monad} is the restriction to the subcategory of
frames and frame homomorphisms. Completeness contributes nothing to
functoriality.

\paragraph{Separation needs only distributivity.}
Theorem~\ref{thm:sharp-iff-point} characterizes what a sharp
valuation \emph{is}, a complement-indicator of a prime ideal, but
not whether there are enough of them to distinguish elements of
$\Omega$. That is a separate fact:

\begin{theorem}[sharp valuations separate points]
\label{thm:sharp-separating}
For every frame $\Omega$ and $a, b \in \Omega$ with $a \not\le b$,
there is a sharp valuation $v$ with $v(a) = 1$ and $v(b) = 0$.
\end{theorem}

\noindent
(Mechanized as \texttt{exists\_sharp\_separating} in
\texttt{Points.lean}.) The proof runs the same prime-ideal
separation theorem that built the counterexample of
Theorem~\ref{thm:hinge}
(\texttt{DistribLattice.prime\_ideal\_of\_disjoint\_filter\_ideal}),
on the principal filter of $a$ and the principal ideal of $b$
rather than on $a \sqcup \neg a$ and $\top$. As a corollary, any two
distinct $a, b \in \Omega$ are distinguished by some sharp
valuation: $\Omega$ need not be linear, so $a \ne b$ alone does not
hand back $a \le b$ or $b \le a$, but it forces at least one failure
of $\le$, and Theorem~\ref{thm:sharp-separating} applies in whichever
direction fails. This is the cash-out of the first qualifying
condition of Section~\ref{sec:conclusion}: the certain valuations do
not merely characterize a single point of the locale in isolation,
together they reconstruct the order on $\Omega$. Reconstructing the
order is not the same as reconstructing the full Heyting structure:
$\neg$ and $\Rightarrow$, as just noted, are adjoints built from an
arbitrary join and carry no meaning over a plain distributive
lattice. What this theorem recovers is the order-theoretic core of
the first qualifying condition, not the negation fragment of
intuitionistic logic, which stays frame-specific.

Tracing the proof of Theorem~\ref{thm:sharp-separating} shows it
never invokes $\bigsqcup$; the only place frame structure enters is
\texttt{DistribLattice.prime\_ideal\_of\_disjoint\_filter\_ideal}
itself, which needs distributivity and nothing beyond it:

\begin{theorem}[separation over distributive lattices]
\label{thm:lvaluation-separating}
For every distributive bounded lattice $\Omega$ and $a, b \in
\Omega$ with $a \not\le b$, there is a sharp lattice valuation $v$
on $\Omega$ with $v(a) = 1$ and $v(b) = 0$.
\end{theorem}

\noindent
(Mechanized as \texttt{LValuation.exists\_sharp\_separating} in
\texttt{GeneralValuation.lean}.) Theorem~\ref{thm:lvaluation-separating}
reproves Theorem~\ref{thm:sharp-separating} with the hypothesis
\texttt{Order.Frame} weakened to \texttt{DistribLattice}, and
nothing else about the argument changes. The condition doing the
work in the sharp/point correspondence is distributivity;
completeness is not load-bearing for either half of it.

\paragraph{Distributivity is not a proof artifact.}
The natural question is whether
Theorem~\ref{thm:lvaluation-separating}'s hypothesis is a genuine
dividing line or an artifact of routing the argument through a
mathlib lemma named for distributive lattices. It is genuine, and
the failure it guards against is total, not merely a failure of
separation:

\begin{definition}[the diamond]
\label{def:m3}
$M_3$ is the five-element lattice with a bottom $z$, a top $o$, and
three pairwise-incomparable atoms $a, b, c$, with $x \sqcup y = o$
and $x \sqcap y = z$ for every two distinct atoms $x, y$.
\end{definition}

\noindent
$M_3$ is the smallest lattice that is not distributive:
$a \sqcap (b \sqcup c) = a \sqcap o = a$, while
$(a \sqcap b) \sqcup (a \sqcap c) = z \sqcup z = z$. It is modular,
so this failure is specific to distributivity rather than some
coarser defect of the order.

\begin{theorem}[no sharp valuation on the diamond]
\label{thm:m3-no-sharp}
No lattice valuation on $M_3$ is sharp.
\end{theorem}

\noindent
(Mechanized as \texttt{M3.no\_sharp\_\allowbreak valuation} in
\texttt{NonDistributive.\allowbreak lean}, together with
\texttt{M3.not\_\allowbreak distributive} for the displayed failure of
distributivity above.) Modularity applied to each of the three
pairs of distinct atoms forces $v(x) + v(y) = 1$, since
$x \sqcup y = \top$ and $x \sqcap y = \bot$. If $v$ were sharp,
$v(a), v(b), v(c) \in \{0, 1\}$ would have to satisfy all three
pairwise sum-to-one equations at once, and a two-element codomain
cannot support that: fixing any one value propagates through the
other two equations to a contradiction with the third. The collapse
on $M_3$ is therefore not a failure of separation, points that
happen to coincide, but a failure of existence: there is no
all-or-nothing valuation on $M_3$ to separate anything with, and the
first qualifying condition of Section~\ref{sec:conclusion} is
unsatisfiable outright.

\paragraph{Where this leaves the boundary.}
Frames are complete distributive lattices.
Theorem~\ref{thm:lvaluation-separating} shows completeness is not
needed for the sharp/point correspondence's separation half, and
Theorem~\ref{thm:m3-no-sharp} shows distributivity cannot be given
up in exchange. The event algebras on which sharp valuations reconstruct the order,
the order-theoretic core of the first qualifying condition of
Section~\ref{sec:conclusion}, are exactly the distributive bounded
lattices, no more and no less. The second qualifying condition and
the negation-dependent machinery of
Sections~\ref{sec:hinge}--\ref{sec:bridges}, slack, the complement
bound, the De Morgan gap, are untouched by this section and remain
tied to frames: nothing here shows, or claims, that $\neg$ or
$\Rightarrow$ survives dropping completeness. This bears on the
orthomodular event lattices of quantum probability discussed in
Section~\ref{sec:related}: giving up distributivity, as those
lattices do, does not merely change which valuations exist, it can
remove sharp valuations altogether, so the construction of this
paper does not, without further modification, have anything to say
about that setting even at the order-theoretic level examined here.
Whether some weaker notion of certainty survives on non-distributive
lattices, in place of the $\{0,1\}$-valued sharp valuations examined
here, is not addressed by this section.
